% Part: first-order-logic
% Chapter: completeness
% Section: downward-ls

\documentclass[../../../include/open-logic-section]{subfiles}

\begin{document}

\olfileid{fol}{com}{dls}
\olsection{Teorema L\"owenheim--Skolem}

Teorema L\"owenheim--Skolem nyatakake yen sawijining teori nduweni
model tanpa wates, mula teori kasebut uga nduweni model kang paling
akeh !!{denumerable}. Akibat langsung saka kasunyatan iki yaiku logika
sepisanan ora bisa nyatakake yen ukuran !!a{structure} iku
!!{nonenumerable}: saben !!{sentence} utawa himpunan !!{sentence}s
kang bisa dicukupi ing kabeh !!{nonenumerable} !!{structure}s uga bisa
dicukupi ing sawijine struktur !!{enumerable}.

\begin{thm}
\ollabel{thm:downward-ls} Yen $\Gamma$ konsisten, himpunan kasebut
nduweni model !!a{enumerable}, yaiku, bisa dicukupi ing
!!a{structure} kang domaine winates utawa !!{denumerable}.
\end{thm}

\begin{proof}
Yen $\Gamma$ konsisten, !!{structure}~$\Struct M$ kang diasilake
dening bukti teorema kelengkapan nduweni domain $\Domain{M}$ kang ora
luwih gedhe tinimbang himpunan term saka basa~$\Lang L$. Mula
$\Struct M$ paling akeh !!{denumerable}.
\end{proof}

\begin{thm}
\ollabel{noidentity-ls} Yen $\Gamma$ minangka himpunan
!!{sentence}s konsisten ing basa logika sepisanan tanpa identitas,
himpunan kasebut nduweni model !!a{denumerable}, yaiku, bisa dicukupi
ing !!a{structure} kang domaine tanpa wates lan !!{enumerable}.
\end{thm}

\begin{proof}
Yen $\Gamma$ konsisten lan ora ngemot ukara kang ngandhut identitas,
!!{structure}~$\Struct M$ kang diasilake dening bukti teorema
kelengkapan nduweni domain $\Domain{M}$ kang padha persis karo
himpunan term saka basa~$\Lang L'$. Mula $\Struct{M}$ iku
!!{denumerable}, amarga $\Trm[L']$ uga mangkono.
\end{proof}

\begin{ex}[Paradoks Skolem]
Teori himpunan Zermelo--Fraenkel~$\Log{ZFC}$ minangka kerangka kang
kuwat banget lan meh kabeh pratelan matematika bisa dinyatakake ing
jerone, kalebu kasunyatan ngenani ukuran himpunan. Contone,
$\Log{ZFC}$ bisa mbuktekake yen himpunan wilangan real~$\Real$ iku
!!{nonenumerable}, bisa mbuktekake Teorema Cantor yen himpunan pangkat
saben himpunan luwih gedhe tinimbang himpunan kasebut dhewe, lan
sapanunggalane. Yen $\Log{ZFC}$ konsisten, kabeh modele tanpa wates,
lan luwih saka iku, kabeh model kasebut ngemot !!{element}s kang
miturut teori iku !!{nonenumerable}, kayata elemen kang nggawe bener
teorema saka~$\Log{ZFC}$ yen himpunan pangkat wilangan asli ana.
Miturut Teorema L\"owenheim--Skolem, $\Log{ZFC}$ uga nduweni model
!!{enumerable}---model kang ngemot himpunan “!!{nonenumerable}”
nanging modele dhewe !!{enumerable}.
\end{ex}

\end{document}

